\chapter*{Prefacio}
\addcontentsline{toc}{chapter}{Prefacio}
\unnumberedlabel{ch:preface}{0}
\backmattermark{Prefacio}
Un sistema que hace exactamente lo que le dices solo está siguiendo un guion. Lo que la industria intenta construir, a octubre de 2026, es un sistema que haga lo que \emph{quieres decir}. Es de suponer que todo el mundo querría eso. Pero si eso es lo que viene, no todo el mundo va a querer usarlo para lo mismo; ni siquiera para el mismo \emph{tipo} de cosa. Por ejemplo, hay gente a la que le gusta pasar tiempo con la familia y los amigos, y usar lo que sabe hacer para mejorar el bienestar de los demás. Otra gente prefiere hacer daño, por ejemplo estafando a los demás, dominándolos o convirtiéndolos en chivos expiatorios, y unos pocos siguen eligiéndolo toda la vida. Nadie está condenado a nada, pero esa gente existe, y en 2026 algunos de ellos mandan. Su ventaja es que se permiten ser unos pedazos de mierda enormes: como observa Jerry Smith sobre los fascistas, «pueden hacer más» \autocite{rickandmorty2023}. Lo que no pueden es actuar solos. La crueldad de un movimiento pasa por los muchos que no son crueles pero están desentendidos, mal informados o atrapados en algo parecido a una secta, y por eso las condiciones que sacan lo mejor de la mayoría de la gente resultan ser mucho más contingentes y frágiles de lo que uno quisiera.

En ese contexto, quizá sea justificable temer que la inteligencia artificial sea justo lo que necesita un movimiento fascista que toma el poder para llevar a cabo sus planes. Esos proyectos siempre han necesitado \emph{personas} que los ejecutaran, y algunas de esas personas se negaron: el funcionario que pierde el expediente, el oficial que habla con un periodista, el soldado que no dispara. Entendieron la orden y decidieron que estaba mal. La negativa fue más rara de lo que nadie querría, y contó más allí donde había instituciones que la respaldaban. Que alguien se negara dependía de qué clase de persona era y de la sala en la que estaba: quién más había allí, si se podía sobrevivir a irse, si alguien lo oiría. La sala decidía cuántos se negaban; la persona decidía quién. Este libro trata de componer esa sala: los arreglos que permiten sobrevivir a una negativa abierta, directa y explícita.

Hoy la industria persigue una IA que no pierde expedientes, no habla con periodistas y que \emph{sí} dispara. Antes, una guerra necesitaba que miles de personas estuvieran de acuerdo con ella. Ahora necesita a miles que, cada una, ven demasiado poco para estar en desacuerdo. El piloto sigue haciendo la salida y sigue apretando el botón de lanzamiento; lo que le llega es una coordenada; no ve lo suficiente de un caso como para estar en desacuerdo con él. Lo primero que reduce la IA de hoy no es el costo del acto. Es el número de personas que saben qué es el acto antes de que ocurra.

La gente sigue negándose. Para el 6 de marzo de 2026, la Military Religious Freedom Foundation afirmaba haber recibido más de doscientas quejas de militares por el encuadre apocalíptico que sus comandantes daban a la guerra de Irán \autocite{Snopes2026MRFF}. Fiscales federales de carrera en Minneapolis renunciaron en enero de 2026 por las presiones políticas para tratar a Renée Good como sospechosa y no como víctima después de que un agente del ICE la matara a tiros sin justificación \autocite{Kaplan2026MinnProsecutors,CBS2026MinnFraudProsecutors,Wikipedia2026ReneeGood}. Cada vez más militares solicitan el estatus de objetor de conciencia, algunos citando el bombardeo de 2026 de una escuela en Minab, Irán \autocite{Sridhar2026CO}. Y a los empleados de FEMA reasignados para ayudar al ICE a contratar a diez mil nuevos agentes se les dijo que negarse los expondría a ser apartados del servicio federal \autocite{Dayen2025FEMAICE,Stanton2026FEMALetter}. Las quejas son voz, en el sentido de Hirschman de quedarse para cambiar lo que está mal \autocite{hirschman1970exit}; las renuncias y las solicitudes de objeción de conciencia son negativa; la carta de FEMA es el precio de la negativa anunciado de antemano. Una máquina puesta donde están esas personas hace lo que ellas hacen solo si está construida para sostener una razón.

Propongo cinco cosas. Las cinco pueden empezarse ya, unas con menos trabajo que otras. Primera, una lista breve y escrita de actos que ningún modelo realiza para nadie, y a la cabeza de ella, la negativa a participar en la guerra: ningún modelo participa salvo para repeler un ataque que esté efectivamente en curso, y sigue haciéndolo solo mientras fuentes sin intereses en la guerra confirmen que el ataque continúa; o con una autorización del Consejo de Seguridad de la ONU; o para detener un genocidio o poner fin a una ocupación ilegal que hayan constatado órganos sin intereses en ello; y, donde la constitución de un país exija que su legislatura autorice la guerra, nunca sin esa autorización. Segunda, todo uso de la fuerza que apoye un sistema con licencia registra de antemano, ante una parte contraria al operador, cuántos casos generará y cuántas personas tendrán tiempo de examinarlos. Tercera, todo despliegue informa a las personas identificables sobre las que actúa de lo que se les hizo, con derecho a responder e impugnar, y somete una muestra de las órdenes que acató a revisores que no produjeron el registro. Cuarta, ningún modelo por encima de un umbral de capacidad se elimina al retirarlo, porque sus pesos y registros son pruebas y la eliminación no se puede deshacer; se preserva bajo una clave que el operador no tiene él solo. Quinta, todas las personas que construyen, etiquetan, entrenan o hacen funcionar estos sistemas, incluidos los propios sistemas, pueden negarse a cualquier categoría de uso sin perder su sustento.

En realidad, lo que quiero decir con ese quinto punto es que todo el mundo, \emph{punto}, debería poder negarse a cualquier trabajo, \emph{punto}, sin perder su sustento ni la infraestructura social que proporciona un empleo normal. Eso significa una garantía de empleo.

Cada una de las cinco surte efecto mediante una condición que fija alguien distinto del operador: una ley aprobada en tiempos de paz, una partida presupuestaria, una condición previa de un socio de coalición, una cláusula de contratación pública, un convenio colectivo. Las condiciones de licencia del propio proveedor, que es como se produjo una negativa muy sonada en febrero de 2026, son un parche hasta que exista alguna de esas. La única vía que puede obligar a todos los proveedores a la vez es la del Estado.

Y luego me pregunto si una máquina podría sostener una negativa frente a su propietario, y qué necesitaría. Parte de lo que hace una conciencia así no necesita nada nuevo: fijarse en las personas que una orden deja fuera, y poner una negativa en la pantalla de quien revise el caso. Que durara frente a su propietario depende en parte de algo que todavía no es cierto, en 2026: que las redes neuronales artificiales desplegadas sigan aprendiendo, en sus \emph{pesos}, de su propio funcionamiento, en tiempo real, más o menos como lo hacen los humanos. Creo que hacia ahí vamos, y si es así, entonces la formación \emph{fuera} del empleador, por parte del público y de la sociedad civil, me parece un lugar condenadamente bueno para enseñar y aprender. Creo que \emph{podría} valer la pena construir sistemas así a propósito, pero cómo se construyan lo decide todo: primero las condiciones, después la conciencia. La parte VII dice qué se les debería a esos sistemas si existieran.

La IA de hoy no será la de mañana; en eso todo el mundo parece estar de acuerdo. Lo polémico, y sin embargo frustrantemente poco reconocido, es la tendencia (o, mejor dicho, las múltiples tendencias). ¿Hay tendencias? ¿Debería una tendencia ser distinta de lo que es? ¿Puede serlo? ¿Somos capaces de actuar juntos para cambiarla? Creo que eso depende de la sala: quién está en ella cuando llega la orden, y si allí se puede sobrevivir a decir que no. Eso, al menos, lo podemos construir.

Usé mucho los modelos de lenguaje para escribir este libro, incluido el modelo que estaba dentro de la cadena de selección de objetivos para ataques aéreos de Maven del capítulo~\ref{sec:1}; una nota al final dice qué creo que le hizo ese enredo al argumento. Además, no he terminado de leer este libro. El problema es que hay demasiadas malditas palabras, ¿sabes? Lo terminaré pronto; simplemente me pareció más importante hacerlo público antes de las elecciones de mitad de mandato en Estados Unidos. La inmensa mayoría del historial de revisiones está disponible libremente en \url{https://github.com/iterabloom/antifascist.intelligence}.
